\chapter{Stage 3: Blow-Up Limits, Kappa-Solutions, and Canonical Neighborhoods}

The analytic controls of the preceding chapter produce the ancient models
seen at high curvature.  Stage 3 classifies those kappa-solutions, proves the
stability and structure of their canonical neighborhoods, and supplies the
canonical model theorem consumed by continuation.  Stage 4 then begins with
the first-failure extension argument and continues through surgery restart,
finite surgery count, and the all-time controlled flow.

\section{Kappa-Solutions and Asymptotic Models}

\begin{theorem}[Two-dimensional kappa-solutions are round]
\label{thm:two-dimensional-kappa-solutions-round}
\dcref{MT:ch9:9.50}
\source{morgan-tian}
\uses{def:kappa-solution,thm:ancient-scalar-curvature-monotonicity}
\notready
Every two-dimensional $\kappa$-solution in the sense of
\ref{def:kappa-solution} is, after a constant parabolic scaling and time
translation, a shrinking compact round surface.  Its orientable double cover
is the shrinking round two-sphere, so the surface is $S^2$ or
$\mathbb{RP}^2$.
\end{theorem}

\begin{proof}
Pass to the orientable double cover when necessary; completeness, bounded
curvature, nonflatness, and noncollapsing persist (with a possibly smaller
constant).  The strong maximum principle makes the Gaussian curvature positive: a zero
would make the ancient solution flat, contrary to nonflatness.  A complete
noncompact positively curved surface is diffeomorphic to the plane.  Its
asymptotic circumference ratio is nonincreasing under the flow; if positive,
a blow-down gives a flat cone and Harnack rigidity makes the solution flat,
while if zero, balls at the curvature scale violate noncollapsing.  Hence the
surface is compact.

On a compact surface, write the normalized metric as
$\widehat g(s)=(-t)^{-1}g(t)$ with $s=-\log(-t)$.  Gauss--Bonnet fixes its
average curvature.  The evolution of the variance of the Gaussian curvature,
together with \ref{thm:ancient-scalar-curvature-monotonicity}, is
nonincreasing and has zero backward limit; otherwise a backward blow-down
would be a nonround compact shrinking soliton, contradicting the integrated
soliton equation on a surface.  Thus the curvature is spatially constant on
every slice of the orientable double cover.  The Ricci-flow equation gives the
shrinking round sphere there; the only free orientation-reversing quotient is
the antipodal quotient $\mathbb{RP}^2$.
\end{proof}

\begin{theorem}[Asymptotic gradient-shrinker production]
\label{thm:asymptotic-gradient-shrinking-soliton}
\dcref{MT:ch9:9.11,MT:ch9:9.42,MT:ch9:9.44}
\source{morgan-tian}
\uses{def:kappa-solution,def:reduced-length,def:reduced-volume,
  lem:reduced-length-support-inequality,lem:minimum-reduced-length-comparison,
  lem:reduced-density-monotonicity,thm:generalized-flow-noncollapsing-interface,
  thm:local-ricci-flow-derivative-estimates,thm:hamilton-compactness-generalized-flows}
\notready
Let $(M,g(t))$ be a complete three-dimensional $\kappa$-solution on
$(-\infty,0]$.  For every basepoint $(p,0)$ there are times
$\tau_j\to\infty$ and points $q_j$ at time $-\tau_j$ with
$l(q_j,\tau_j)\le3/2$ such that the parabolically rescaled pointed flows
based at $(q_j,-1)$ subconverge smoothly on compact spacetime sets to a
complete nonflat ancient gradient shrinking Ricci soliton with bounded
nonnegative curvature and the inherited noncollapsing estimate.
\end{theorem}
\begin{proof}
This is the asymptotic-model production theorem proved in Morgan--Tian,
Chapter 9, Theorems 9.11, 9.42, and 9.44.  The displayed hypotheses are the
ones used there: minimum reduced length chooses the points, the support and
local derivative bounds give compactness, and reduced-density monotonicity
gives the shrinking-soliton equality.  The generalized Hamilton compactness
interface in the dependency list supplies the smooth pointed limit.  No
classification conclusion is included in this imported production contract.
\end{proof}

\begin{theorem}[Three-dimensional gradient-shrinker classification]
\label{thm:three-dimensional-gradient-shrinker-classification}
\dcref{MT:ch9:9.53,MT:ch9:9.54}
\source{morgan-tian}
\uses{def:kappa-solution,thm:curvature-null-splitting,
  thm:two-dimensional-kappa-solutions-round}
\notready
Let $(N,h(t),f)$ be a complete nonflat three-dimensional ancient gradient
shrinking soliton with bounded nonnegative curvature and positive
$\kappa$-noncollapsing on every finite scale.  Then it is either a compact
round space form or one of
\[
 S^2\times\mathbb R,\qquad
 \mathbb{RP}^2\times\mathbb R,\qquad
 \mathbb{RP}^2\widetilde\times\mathbb R,
\]
with their standard shrinking product metrics.  Every noncompact member is
one of the three displayed cylindrical models.
\end{theorem}
\begin{proof}
This is the exact imported classification in Morgan--Tian, Chapter 9,
Theorems 9.53--9.54.  Its hypotheses include completeness, bounded
nonnegative curvature, nonflatness, and finite-scale noncollapsing.  The
null-splitting and two-dimensional roundness nodes record the source's split
case; the remaining soliton classification is imported with its precise
source location rather than replaced by an unstated maximum-principle step.
\end{proof}

\begin{theorem}[Asymptotic models of three-dimensional kappa-solutions]
\label{thm:kappa-solution-asymptotic-models}
\dcref{MT:ch9:9.11,MT:ch9:9.42,MT:ch9:9.44,MT:ch9:9.53,MT:ch9:9.54}
\source{morgan-tian}
\notready
Let $(M,g(t))$ be a three-dimensional $\kappa$-solution and fix $(p,0)$.
There are $\tau_j\to\infty$ and points $q_j$ at time $-\tau_j$ with reduced
length at most $3/2$ such that the parabolically rescaled pointed flows based
at $(q_j,-1)$ converge to a complete nonflat gradient shrinking soliton with
bounded nonnegative curvature.  The limit is a compact round space form or
one of the three cylindrical models
\[
 S^2\times\mathbb R,\qquad
 \mathbb{RP}^2\times\mathbb R,\qquad
 \mathbb{RP}^2\widetilde\times\mathbb R,
\]
where the last space is the twisted line bundle.  If the original solution
is nonround, a sequence can be chosen with one of these three noncompact
limits.
\end{theorem}

\begin{proof}
\uses{thm:asymptotic-gradient-shrinking-soliton,
  thm:three-dimensional-gradient-shrinker-classification,
  thm:ancient-scalar-curvature-monotonicity}
The production interface \ref{thm:asymptotic-gradient-shrinking-soliton}
supplies the pointed soliton limit and the reduced-length normalization.
Apply the imported classification
\ref{thm:three-dimensional-gradient-shrinker-classification} to that limit.
The remaining source argument uses reduced-volume equality to distinguish the
round case from a noncompact asymptotic sequence; this is the precise
nonround alternative asserted below.
\end{proof}

\begin{proposition}[Universal noncollapsing for nonround kappa-solutions]
\label{prop:universal-nonround-kappa}
\dcref{MT:ch9:9.58}
\source{morgan-tian}
\uses{def:kappa-solution}
\notready
There is a universal $\kappa_*>0$ such that every nonround
three-dimensional $\kappa$-solution is uniformly $\kappa_*$-noncollapsed:
for every $t\le0$, $x\in M_t$, and every $r>0$ for which the backward
parabolic ball is defined and satisfies $|\operatorname{Rm}|\le r^{-2}$,
\[
 \operatorname{Vol}_{g(t)}B_{g(t)}(x,r)\ge\kappa_*r^3.
\]
The constant $\kappa_*$ depends only on the dimension, not on the supplied
$\kappa$ or on the particular nonround solution; equivalently, every such
flow is a $\kappa_*$-solution in the sense of \ref{def:kappa-solution}.
\end{proposition}

\begin{proof}
\uses{thm:kappa-solution-asymptotic-models,
  thm:three-dimensional-gradient-shrinker-classification,
  thm:generalized-flow-noncollapsing-interface,
  lem:reduced-volume-noncollapsing-transfer}
The universal quantifier is the content of Morgan--Tian
Proposition~9.58, not a consequence of one selected asymptotic point.  Its
contradiction argument starts with a violating triple $(t,x,r)$, translates
$t$ to zero, and rescales by $r^{-2}$; the curvature and backward-domain
hypotheses are preserved.  The asymptotic-model theorem
\ref{thm:kappa-solution-asymptotic-models} and the finite classification
\ref{thm:three-dimensional-gradient-shrinker-classification} reduce the
violating sequence to one of the three cylindrical models.  Choose a model
time $\bar\tau$, a reduced-length bound $l_0$, and a compact model ball with
reduced-volume contribution $V_*>0$; the finite model list makes these
constants universal.  Smooth convergence gives $V_*/2$ on the rescaled
sequence, and
\ref{lem:reduced-volume-noncollapsing-transfer} converts it to a volume
ratio bounded below by a fixed $\kappa_*>0$, contradicting the violating
triple.  Since the initial triple was arbitrary, undoing the rescaling gives
the displayed estimate at every slice, center, and admissible finite scale.
\end{proof}

\begin{theorem}[Normalized kappa-solution compactness]
\label{thm:kappa-solution-compactness}
\dcref{MT:ch9:9.59,MT:ch9:9.60,MT:ch9:9.62,MT:ch9:9.63,MT:ch9:9.64,MT:ch9:9.65,MT:ch9:9.66,MT:ch9:9.67,MT:ch9:9.68,MT:ch9:9.69,MT:ch9:9.70,MT:ch9:9.71}
\source{morgan-tian}
\notready
Fix $\kappa>0$.  Let $(M_j,g_j(t),p_j)$ be based three-dimensional
$\kappa$-solutions for this common value of $\kappa$, with
$R(p_j,0)=1$.  A subsequence converges smoothly on compact spacetime
subsets to a based $\kappa$-solution
\[
 (M_\infty,g_\infty(t),p_\infty)
\]
with $R(p_\infty,0)=1$.
In particular, for every $A<\infty$ and $m\ge0$ there is $C(A,m)<\infty$
such that
\[
 \sup_{B_{g_j(0)}(p_j,A)\times[-A,0]}
 |\nabla^m\operatorname{Rm}_{g_j}|\le C(A,m).
\]
\end{theorem}

\begin{proof}
\uses{def:kappa-solution,thm:hamilton-compactness-generalized-flows,thm:local-ricci-flow-derivative-estimates,thm:local-volume-injectivity-radius-control,thm:curvature-null-splitting,thm:ancient-scalar-curvature-monotonicity,lem:nonflat-terminal-cone-exclusion}
First obtain a uniform scalar-curvature bound on every fixed zero-time ball.
If this failed, choose a nearest point $q_j$ at which the curvature scale is
smaller than its distance scale from $p_j$, and rescale at $q_j$.  The
point-picking choice bounds curvature on a ball whose rescaled radius tends
to infinity.  If the volume ratio of these balls stayed uniformly positive,
$\kappa$-noncollapsing and Hamilton compactness would give a complete
nonflat limit with a zero-curvature point, contradicting
\ref{thm:curvature-null-splitting}.  Otherwise choose the first radius at
which the volume ratio is
half Euclidean.  Volume comparison and the point-picking curvature bound
give a complete flat limit whose unit ball has half the Euclidean volume.
It cannot be Euclidean, while every nontrivial complete flat quotient has
zero asymptotic volume ratio and hence is not noncollapsed on all scales.
This contradiction proves the required based-ball curvature bound.

By \ref{thm:ancient-scalar-curvature-monotonicity}, the same bound holds on
the corresponding backward cylinders.  The local estimates
\ref{thm:local-ricci-flow-derivative-estimates} bound every covariant
derivative on smaller cylinders.
Noncollapsing supplies the base-ball volume lower bound in
\ref{thm:hamilton-compactness-generalized-flows}; diagonalization in the
radius, derivative order, and backward time therefore gives an ancient
smooth limit complete at time zero.  Its nonnegative curvature makes every
earlier slice complete.  It is nonflat because its base scalar curvature is
one, has nonnegative curvature by smooth convergence, and remains
$\kappa$-noncollapsed.

If a limit time slice had unbounded curvature, repeat the point selection
farther and farther from the basepoint.  Blowing down the separating annuli
would produce a nonflat metric cone realized on a finite time slice (translate
that slice to time $0$),
contrary to \ref{lem:nonflat-terminal-cone-exclusion}.  Thus every limit
slice has bounded curvature, so the limit satisfies
\ref{def:kappa-solution}.  The derivative bounds obtained above are uniform
before passage to the subsequence and give the last assertion.
\end{proof}

\section{Kappa-Solution Geometry and Stability}

\begin{lemma}[Stability of canonical neighborhoods]
\label{lem:canonical-neighborhood-stability}
\dcref{MT:ch9:9.79,MT:ch9:9.80}
\source{morgan-tian}
\notready
Fix $\epsilon<1/4$ and $C<\infty$.
\begin{enumerate}
\item A $(C,\epsilon)$-cap, controlled component, or round component
      persists under sufficiently small
      $C^{\lfloor\epsilon^{-1}\rfloor}$ perturbations of its metric and
      under a sufficiently small change of center.
\item An evolving $\epsilon$-neck of normalized backward length strictly
      greater than one persists as a strong $\epsilon$-neck under
      sufficiently small smooth perturbations of the flow.
\item If $C'>C$, $\epsilon'>\epsilon$, and the center satisfies $R(x)\le2$,
      then there is $\delta>0$ such that every metric which is $\delta$-close
      in $C^{\lfloor\epsilon^{-1}\rfloor}$ to the given
      $(C,\epsilon)$-canonical neighborhood contains a
      $(C',\epsilon')$-canonical neighborhood of $x$.
\end{enumerate}
In particular, strong neck centers form an open spacetime set.
\end{lemma}

\begin{proof}
\uses{def:strong-canonical-neighborhood-control,def:canonical-cap,def:epsilon-neck-and-scale}
Curvature ratios, normalized diameters and volumes, core-volume lower
bounds, and curvature-derivative quantities in the definitions all vary
continuously in the stated topology.  A neck chart survives a small metric
perturbation; precomposing it with a small axial translation recenters the
chart.  For a cap, first keep the compact core and its strict estimates,
then truncate and recenter the neck end.  These observations prove the
first assertion and the metric part of the third.

For an evolving neck, the excess normalized time beyond one absorbs the
small change in curvature scale caused by perturbation.  Restricting the
persisting cylinder to normalized length one gives the required strong
neck.  Enlarging $C$ and $\epsilon$ turns all weak boundary inequalities
into strict ones, proving the last metric assertion.  The same argument on
a slightly longer spacetime cylinder proves openness of strong neck centers.
\end{proof}

\begin{proposition}[Curvature trichotomy for kappa-solutions]
\label{prop:kappa-solution-curvature-trichotomy}
\dcref{MT:ch9:9.83}
\source{morgan-tian}
\notready
Every three-dimensional $\kappa$-solution is exactly one of the following
pairwise disjoint alternatives:
\begin{enumerate}
\item its curvature operator is positive definite (equivalently, every
      sectional curvature is positive) on every time-slice;
\item after a constant parabolic scaling and time translation, it is the
      shrinking round product $S^2\times\mathbb R$;
\item after such a normalization, it is the shrinking round flow on the
      total space of a real line bundle over $\mathbb{RP}^2$, whose twofold
      cover is $S^2\times\mathbb R$.  This includes the trivial product
      $\mathbb{RP}^2\times\mathbb R$ and the twisted line bundle
      $\mathbb{RP}^2\widetilde\times\mathbb R$.
\end{enumerate}
\end{proposition}

\begin{proof}
\uses{def:kappa-solution,thm:curvature-null-splitting,thm:two-dimensional-kappa-solutions-round}
If the curvature operator has a zero eigenvalue on some time-slice, its null
distribution is parallel by \ref{thm:curvature-null-splitting}.  The universal
cover consequently splits for all earlier times as a two-dimensional ancient
solution times a line or a circle.  A circle factor has fixed length while the
curved factor expands backward, so curvature-scale balls collapse as
$t\to-\infty$, contradicting $\kappa$-noncollapsing.  The flat factor is
therefore a line.  The two-dimensional factor is the shrinking round sphere
by \ref{thm:two-dimensional-kappa-solutions-round}; the possible free
line-bundle quotients are exactly alternatives (2) and (3).  If no zero
eigenvalue occurs, nonnegative curvature operator makes every sectional
curvature positive, giving (1), so the branches are disjoint.
\end{proof}

\begin{theorem}[Soul structure in nonnegative curvature]
\label{thm:soul-structure-nonnegative-curvature}
\dcref{MT:ch2:2.7}
\source{morgan-tian}
\notready
Every connected complete noncompact manifold of nonnegative sectional
curvature contains a compact totally geodesic, totally convex soul and is
diffeomorphic to its normal bundle.  If the sectional curvature is positive,
the soul is a point.  The proof uses the Cheeger--Gromoll soul theorem,
smooth Schoenflies/annulus separation for the resulting embedded spheres, and
the covering-space lifting theorem for the normal exponential map; these are
part of the imported topology contract and are not inferred from curvature
alone.
\end{theorem}

\begin{proof}
The existence of a compact totally convex soul and the diffeomorphism by the
normal exponential map are the Cheeger--Gromoll soul theorem, in the exact
complete/nonnegative-curvature form invoked by Morgan--Tian, Chapter 2,
Theorem 2.7.  The positive-curvature point-soul conclusion is the corollary
used there: a positive sectional curvature soul cannot contain a nonconstant
geodesic.  We import these two statements with their complete-manifold and
sectional-curvature hypotheses; no finite-exhaustion or normal-exponential
argument is being treated as a proof of the imported theorem here.
\end{proof}

\begin{lemma}[Imported interface: a thin neck separates a point soul from the end]
\label{lem:positive-curvature-neck-separation}
\dcref{MT:ch2:2.20}
\uses{def:epsilon-neck-and-scale,thm:soul-structure-nonnegative-curvature}
\source{morgan-tian}
\notready
For sufficiently small $\epsilon$, let $M^3$ be complete, noncompact, and
positively curved with point soul $p$.  The central sphere of every
$\epsilon$-neck disjoint from $p$ separates $p$ from the unique end.  Two
disjoint $\epsilon$-necks which avoid $p$ have disjoint central spheres that
bound a region diffeomorphic to $S^2\times[0,1]$.
\end{lemma}

\begin{proof}
The point-soul conclusion in
\ref{thm:soul-structure-nonnegative-curvature} identifies $M$ with the normal
bundle of $p$, hence with $\mathbb R^3$, and its distance spheres separate the
soul from the unique end.  In an $\epsilon$-neck disjoint from $p$, radial
normal geodesics meet the central cross-section transversely: the neck metric
is $C^2$-close to the round cylinder, while positive sectional curvature
prevents a radial geodesic from returning to the soul side.  Nearest-point
projection therefore carries the central sphere by an isotopy to a distance
level sphere, proving separation.  For two disjoint necks, order their
cross-sections by this radial coordinate.  The region between them has no
critical point of the distance function, so the normal exponential map gives
an $S^2\times[0,1]$ product.  This is the comparison and isotopy argument of
Morgan--Tian Lemma 2.20.
\end{proof}

\begin{proposition}[Positive curvature excludes arbitrarily small necks]
\label{prop:positive-curvature-neck-scale-lower-bound}
\dcref{MT:ch2:2.19}
\source{morgan-tian}
\notready
For every sufficiently small $\epsilon>0$, a complete positively curved
Riemannian three-manifold contains no sequence of $\epsilon$-necks whose
scales tend to zero.
\end{proposition}

\begin{proof}
\uses{def:epsilon-neck-and-scale,thm:soul-structure-nonnegative-curvature,thm:busemann-function-superharmonicity,lem:epsilon-neck-geodesic-axis-alignment,lem:positive-curvature-neck-separation}
The assertion is immediate on a compact manifold.  Otherwise
\ref{thm:soul-structure-nonnegative-curvature} gives a point soul $p$ and
identifies the manifold with $\mathbb R^3$.  By
\ref{lem:positive-curvature-neck-separation}, disjoint necks away from the
soul can be ordered toward the end.

Choose two such ordered necks $N_1,N_2$, with centers $x_i$ and scalar
curvatures $R_i$.  Let $B$ be the Busemann function of a ray from $p$, and
choose a nonnegative cutoff $\psi$ which is one between the central regions
and changes only along the middle thirds of the two necks.  The alignment
estimate \ref{lem:epsilon-neck-geodesic-axis-alignment} compares $\nabla B$
with the signed axial direction on both transition regions.  Weak
superharmonicity from \ref{thm:busemann-function-superharmonicity} and
integration by parts give
\[
 0\le \int_M\langle\nabla B,\nabla\psi\rangle\,d\mu
   =\bigl(\alpha_2R_2^{-1}-\alpha_1R_1^{-1}\bigr)
      \operatorname{Vol}(S^2,h_0),
\]
where the cylindrical neck estimates give
$|\alpha_i-1|\le c\epsilon$.  After decreasing $\epsilon$ so that
$c\epsilon<1/3$, this implies $R_2\le2R_1$.  Fixing the first neck therefore
bounds below the scale $R^{-1/2}$ of every later neck.  On the compact region
on the soul side, scalar curvature is bounded, so neck scales have a positive
lower bound there as well.  Hence no sequence of neck scales can tend to
zero.
\end{proof}

\begin{lemma}[Point-picking and the splitting-limit producer]
\label{lem:kappa-point-picking-line}
\dcref{MT:ch9:9.39}
\source{morgan-tian}
\uses{def:kappa-solution,thm:ancient-scalar-curvature-monotonicity,
  thm:hamilton-compactness-generalized-flows,
  thm:busemann-function-superharmonicity,thm:curvature-null-splitting}
\notready
Let $(M,g(t))$, $-\infty<t\le0$, be a complete noncompact three-dimensional
$\kappa$-solution and let $p_j\to\infty$ satisfy
$R(p_j,-1)d(p,p_j)^2\to\infty$.  The scalar-curvature point-picking rule on a
ball $B_{g(-1)}(p_j,d_j/2)$, where $d_j=d_{g(-1)}(p,p_j)$, produces
$q_j$ and $Q_j=R(q_j,-1)$ such that
\[
 q_j\in B(p_j,d_j/2),\qquad Q_j\ge R(p_j,-1),\qquad
 Q_jd(p,q_j)^2\longrightarrow\infty,
\]
and
\[
 R(\,\cdot\,,t)\le4Q_j
 \quad\hbox{on}\quad
 B_{g(-1)}\!\left(q_j,
   \frac{d_jR(p_j,-1)^{1/2}}{4Q_j^{1/2}}\right)
 \quad(t\le-1).
\]
The radii of these balls tend to infinity after scaling by $Q_j$.  A
subsequence of the pointed rescaled flows therefore converges smoothly to a
complete ancient product flow $(N,h(t))\times(\mathbb R,ds^2)$ with $N$
nonflat and of bounded nonnegative curvature.
\end{lemma}

\begin{proof}
Apply the compact-ball point-picking procedure to $R^{1/2}$ on
$B(p_j,d_j/2)$.  Starting at $p_j$, if the factor-two bound for $R^{1/2}$
fails on the prescribed smaller ball, move to a point where it doubles and
halve the allowed moving radius.  The successive moving distances have sum
less than $d_j/2$, while the curvature values double.  Compactness of the
closed ball and boundedness of $R$ force the procedure to stop.  Its terminal
point $q_j$ has the displayed factor-four scalar bound;
because $d(p,q_j)\ge d_j/2$ and $Q_j\ge R(p_j,-1)$, both asserted divergent
products follow.  Hamilton's Harnack inequality, recorded in
\ref{thm:ancient-scalar-curvature-monotonicity}, extends the factor-four bound
to every $t\le-1$.  Noncollapsing and Hamilton compactness now give the ancient
geometric limit on balls exhausting the limit slice.

The line is supplied by Morgan--Tian's splitting-at-infinity argument, not by
the point-picking procedure alone.  In the metrics $Q_jg(-1)$,
minimizing arcs from $q_j$ to $p$ have lengths tending to infinity.  Choose a
farther $q_{m(j)}$ whose minimizing direction at $p$ is arbitrarily close to
that of $q_j$ and whose distance from $p$ is arbitrarily larger.  Toponogov's
hinge comparison shows that the limiting ray toward $p$ and the limiting ray
toward $q_{m(j)}$ meet at angle $\pi$ and form a complete minimizing line.

To obtain the product without another hidden interface, let $b_+$ and $b_-$
be the two Busemann functions of this line.  The triangle inequality gives
$b_++b_-\ge0$ with equality on the line, while
\ref{thm:busemann-function-superharmonicity} makes the sum weakly
superharmonic.  The strong maximum principle gives $b_++b_-=0$, so $b_+$ is
smooth and harmonic.  Its asymptotic rays give $|\nabla b_+|=1$; the Bochner
identity then gives $\operatorname{Hess}b_+=0$.  The complete parallel flow of
$\nabla b_+$ identifies the slice isometrically with
$b_+^{-1}(0)\times\mathbb R$.  Finally,
\ref{thm:curvature-null-splitting} propagates this product through the ancient
flow.  The basepoint normalization gives nonflatness of $N$, and the
factor-four estimate gives bounded curvature.
\end{proof}

\begin{theorem}[Splitting at infinity for kappa-solutions]
\label{thm:kappa-solution-splitting-at-infinity}
\dcref{MT:ch9:9.39,MT:ch9:9.52}
\source{morgan-tian}
\notready
Let $(M,g(t))$, $-\infty<t\le0$, be a noncompact three-dimensional
$\kappa$-solution.  If $p_j\to\infty$ and
\[
 R(p_j,-1)d_{g(-1)}(p,p_j)^2\longrightarrow\infty,
\]
then there are points $q_j\to\infty$, scales $Q_j=R(q_j,-1)$, and a
subsequence for which the pointed rescaled flows converge smoothly to
\[
 (\Sigma,h(t))\times(\mathbb R,ds^2),
\]
where $(\Sigma,h(t))$ is a shrinking compact round surface.  Thus
$\Sigma$ is $S^2$ or $\mathbb{RP}^2$; if $M$ is orientable, it is $S^2$.
\end{theorem}

\begin{proof}
\uses{def:kappa-solution,lem:kappa-point-picking-line,
  thm:two-dimensional-kappa-solutions-round}
Apply \ref{lem:kappa-point-picking-line}.  Its surface factor inherits
completeness, nonflatness, bounded nonnegative curvature and a quantitative
noncollapsing constant from the product limit, so its orientable double cover
is a two-dimensional $\kappa'$-solution.  By
\ref{thm:two-dimensional-kappa-solutions-round} that cover is a shrinking
round $S^2$.  The only fixed-point-free quotient of a round $S^2$ is the
antipodal quotient, so the surface factor is $S^2$ or $\mathbb{RP}^2$ with its
round shrinking metric.  Smooth pointed convergence preserves orientation;
hence the second alternative is absent when $M$ is orientable.
\end{proof}

\begin{theorem}[Structure of noncompact kappa-solutions]
\label{thm:noncompact-kappa-solution-structure}
\dcref{MT:ch9:9.84,MT:ch9:9.85,MT:ch9:9.87,MT:ch9:9.88}
\source{morgan-tian}
\notready
For every sufficiently small $\epsilon>0$ there is
$C_0=C_0(\epsilon)<\infty$ such that the time-zero slice of every
noncompact three-dimensional $\kappa$-solution containing no embedded
two-sided $\mathbb{RP}^2$ is either a strong $\epsilon$-tube or a
$(C_0,\epsilon)$-cap followed by a strong $\epsilon$-tube.  The cap core has
uniform scale-normalized curvature, derivative, diameter, and volume bounds.
\end{theorem}

\begin{proof}
\uses{def:kappa-solution,prop:universal-nonround-kappa,prop:kappa-solution-curvature-trichotomy,thm:soul-structure-nonnegative-curvature,thm:kappa-solution-splitting-at-infinity,thm:kappa-solution-compactness,lem:canonical-neighborhood-stability,def:canonical-neck-cap-regions,thm:canonical-neck-cap-topology}
Consider first the positive-curvature case and normalize at a soul $p$ by
$R(p,0)=1$, using \ref{thm:soul-structure-nonnegative-curvature}.  Suppose
points $q_j$ escaping every fixed ball were not strong
$\epsilon$-neck centers.  If
$d(p,q_j)^2R(q_j,0)\to\infty$,
\ref{thm:kappa-solution-splitting-at-infinity} gives the shrinking round
cylinder, and
\ref{lem:canonical-neighborhood-stability} makes $q_j$ a neck center, a
contradiction.  If those products remain bounded, then
$R(q_j,0)\to0$ while the unit-curvature soul stays at bounded distance in
the $R(q_j,0)$-rescaling.  The based curvature bound in
\ref{thm:kappa-solution-compactness} rules this out.

The universal constant in \ref{prop:universal-nonround-kappa} makes the
radius beyond which all points are neck centers uniform.  Otherwise
normalize a sequence of souls and take the first non-neck point beyond an
ever longer neck annulus.  Compactness at the souls transfers a fixed neck
annulus to the sequence.  At the outer first-failure point the preceding
scale-separation argument again gives a cylindrical limit, contradicting
stability.  Compactness on the remaining soul ball supplies uniform
curvature ratios, derivative bounds, volume bounds, and diameter bounds.
The neck spheres separate the unique end, and
\ref{thm:canonical-neck-cap-topology} identifies the bounded remainder as a
cap core.

In the zero-curvature cases of
\ref{prop:kappa-solution-curvature-trichotomy}, the round cylinder is already
a strong tube.  Its nontrivial orientable quotient has one controlled cap
core and a cylindrical end.  The excluded two-sided
$\mathbb{RP}^2$ removes the remaining exceptional product.  Increasing
$C_0$ to dominate the uniform core estimates proves the theorem.
\end{proof}

\begin{theorem}[Structure of compact kappa-solutions]
\label{thm:compact-kappa-solution-structure}
\dcref{MT:ch9:9.89,MT:ch9:9.90,MT:ch9:9.91}
\source{morgan-tian}
\notready
For every sufficiently small $\epsilon>0$ there is
$C_1=C_1(\epsilon)<\infty$ such that a compact three-dimensional
$\kappa$-solution is round, is a positively curved $C_1$-component, or is
the union of two disjoint $(C_1,\epsilon)$-cap cores joined by a strong
$\epsilon$-tube.  In the nonround cases its time-zero slice is
diffeomorphic to $S^3$ or $\mathbb{RP}^3$.
\end{theorem}

\begin{proof}
\uses{def:kappa-solution,prop:universal-nonround-kappa,prop:kappa-solution-curvature-trichotomy,thm:kappa-solution-compactness,thm:noncompact-kappa-solution-structure,lem:canonical-neighborhood-stability,thm:canonical-neck-cap-topology}
The splitting alternatives in
\ref{prop:kappa-solution-curvature-trichotomy} are noncompact, so a compact
nonround solution has positive sectional curvature and finite fundamental
group.  Normalize its maximum scalar curvature to one.  If a sequence has
unbounded normalized diameter and contains a point with neither a cap nor a
neck neighborhood, compactness based at that point gives a noncompact
$\kappa$-solution.  The point lies in the cap/tube decomposition of
\ref{thm:noncompact-kappa-solution-structure}, and stability transfers that
neighborhood back, a contradiction.  Hence every large-diameter solution is
a double-capped strong tube.  The topology classification in
\ref{thm:canonical-neck-cap-topology}, together with finite fundamental
group, leaves $S^3$ and $\mathbb{RP}^3$.

For bounded normalized diameter,
\ref{prop:universal-nonround-kappa} puts every nonround member in one fixed
noncollapsed family, which is compact by
\ref{thm:kappa-solution-compactness}.  A nonround limit retains positive
curvature, so sectional-curvature ratios, diameter, volume, and curvature
derivative ratios have positive finite extrema over this family.  Taking
$C_1$ larger than those extrema makes the whole slice a controlled
component.  A round limit gives the round alternative.  These cases exhaust
the compact family and prove the assertion.
\end{proof}

\begin{theorem}[Uniform structure of kappa-solutions]
\label{thm:kappa-solution-structure-summary}
\dcref{MT:ch9:9.93}
\source{morgan-tian}
\notready
For every sufficiently small $\epsilon>0$ there is
$C=C(\epsilon)<\infty$ such that every three-dimensional
$\kappa$-solution is exactly one of the following seven source alternatives.
For the compact nonround branches, put
$R_{\max}=\max_{M}R(\cdot,0)$ and fix the diameter threshold
$D_*=D_*(\epsilon)$ supplied by the large-diameter structure theorem.  The
round branch is disjoint from the nonround branches, and the two compact
nonround branches are separated by this threshold:
\begin{enumerate}
\item the flow is round for all $t\le0$, and its time-zero slice is
      $S^3/\Gamma$ for a finite subgroup $\Gamma<SO(4)$ acting freely;
\item (nonround) the time-zero slice is compact and positively curved, has
      $\operatorname{diam}M\le D_*R_{\max}^{-1/2}$ and
      $C^{-1}R(x,0)^{-1/2}<\operatorname{diam}M<C R(x,0)^{-1/2}$,
      $C^{-1}R(x,0)^{-3/2}<\operatorname{Vol}M<C R(x,0)^{-3/2}$,
      and $C^{-1}R(x,0)<K(P_y)<C R(x,0)$ for every $y$ and $2$-plane
      $P_y$; it is diffeomorphic to $S^3$ or $\mathbb{RP}^3$;
\item (nonround) the time-zero slice is positively curved and is a double
      $C$-capped strong $\epsilon$-tube with
      $\operatorname{diam}M>D_*R_{\max}^{-1/2}$, hence is diffeomorphic to
      $S^3$ or $\mathbb{RP}^3$;
\item (nonround) the time-zero slice is positively curved and is a $C$-capped strong
      $\epsilon$-tube diffeomorphic to $\mathbb R^3$;
\item the time-zero slice is the quotient of a round $S^2\times\mathbb R$
      by a free orientation-preserving involution, is a $C$-capped strong
      $\epsilon$-tube, and is diffeomorphic to a punctured $\mathbb{RP}^3$;
\item the time-zero slice is the product of a round $S^2$ and $\mathbb R$
      and is a strong $\epsilon$-tube;
\item the time-zero slice is the product $\mathbb{RP}^2\times\mathbb R$
      with the constant-curvature metric on $\mathbb{RP}^2$.
\end{enumerate}
In alternatives (3)--(5), points are in cap cores or are centers of strong
$\epsilon$-necks; in alternatives (1), (2), and (7) the source gives the
corresponding whole-component description.  Uniformly on the flow,
\[
 \sup_{p\in M,\,t\le0}\frac{|\nabla R(p,t)|}{R(p,t)^{3/2}}\le C,
 \qquad
  \sup_{p\in M,\,t\le0}\frac{|\partial_tR(p,t)|}{R(p,t)^2}\le C.
\]
\end{theorem}

\begin{proof}
\uses{prop:universal-nonround-kappa,prop:kappa-solution-curvature-trichotomy,
  thm:noncompact-kappa-solution-structure,thm:compact-kappa-solution-structure,
  thm:kappa-solution-compactness,thm:local-ricci-flow-derivative-estimates,
  thm:curvature-null-splitting}
Apply \ref{prop:kappa-solution-curvature-trichotomy}.  The round branch is
alternative (1).  If the curvature has a null direction, the finite-interval
splitting theorem \ref{thm:curvature-null-splitting}, followed by the
two-dimensional roundness theorem, gives precisely the twisted quotient, the
round $S^2\times\mathbb R$, or the round
$\mathbb{RP}^2\times\mathbb R$, namely alternatives (5)--(7).  In the
strictly positive branch, the noncompact structure theorem gives alternative
(4), while the compact structure theorem gives alternatives (2) or (3).
These source alternatives are disjoint and exhaust the trichotomy.  The
topology clauses in those source contracts give the stated diffeomorphism
types and no additional case.
Choose $C$ larger than all constants in the two structure theorems and in
\ref{thm:kappa-solution-compactness}; the compactness and universal
noncollapsing inputs make the normalized derivative ratios uniform.  Scaling
back gives both displayed supremum bounds.
\end{proof}

\section{Canonical Neighborhood Models}

\begin{theorem}[Canonical neighborhoods in kappa-solutions]
\label{thm:kappa-solution-canonical-neighborhoods}
\dcref{MT:ch9:9.94,MT:ch9:9.95}
\source{morgan-tian}
\notready
For every sufficiently small $\epsilon>0$ there is
$C=C(\epsilon)<\infty$ such that every point of every three-dimensional
$\kappa$-solution has a strong $(C,\epsilon)$-canonical neighborhood,
unless the solution is the round product $\mathbb{RP}^2\times\mathbb R$.
Moreover, if generalized flows with no embedded two-sided
$\mathbb{RP}^2$ converge smoothly to a $\kappa$-solution, then their
basepoints have strong $(C,\epsilon)$-canonical neighborhoods for all
sufficiently large indices.
\end{theorem}

\begin{proof}
\uses{def:strong-canonical-neighborhood-control,thm:kappa-solution-structure-summary,lem:canonical-neighborhood-stability}
Every nonexceptional alternative in
\ref{thm:kappa-solution-structure-summary} is one of the four alternatives
in \ref{def:strong-canonical-neighborhood-control}; taking the same uniform
$C(\epsilon)$ proves existence.  The only remaining alternative is the round
$\mathbb{RP}^2\times\mathbb R$ product.

For the convergence assertion, the limiting basepoint has a canonical
neighborhood by the first part.  A cap, controlled component, or round
component transfers directly by the first part of
\ref{lem:canonical-neighborhood-stability}.  A strong neck extends to a
slightly longer normalized backward interval, so the second part of that
lemma transfers it as a strong neck.  The limit cannot be the exceptional
product: an embedded central $\mathbb{RP}^2$ has trivial normal bundle and
persists under smooth convergence, contrary to the hypothesis on the
approximating slices.  Thus the basepoints eventually have the asserted
canonical neighborhoods.
\end{proof}

\begin{theorem}[Canonical neighborhoods in the standard solution]
\label{thm:standard-solution-canonical-neighborhoods}
\dcref{MT:ch12:12.28,MT:ch12:12.32,MT:ch12:12.36}
\source{kleiner-lott}
\notready
For every sufficiently small $\epsilon>0$ there is
$C_{\mathrm{std}}(\epsilon)<\infty$ such that every point $(q,t)$ of the
standard solution has one of the following neighborhoods:
\begin{enumerate}
\item a $(C_{\mathrm{std}},\epsilon)$-cap whose core contains $q$;
\item an evolving $\epsilon$-neck reaching time zero, with its initial slice
      disjoint from the cap core;
\item an evolving $\epsilon$-neck defined for normalized backward time
      $1+\epsilon$.
\end{enumerate}
The alternatives persist for flows converging smoothly to the standard
solution on the corresponding parabolic neighborhoods.
\end{theorem}

\begin{proof}
\uses{def:strong-canonical-neighborhood-control,def:standard-surgery-cap-model,thm:hamilton-compactness-generalized-flows,thm:kappa-solution-canonical-neighborhoods,lem:canonical-neighborhood-stability}
The standard solution is rotationally symmetric, positively curved, and
asymptotic to the evolving round cylinder.  Its scalar curvature tends
uniformly to infinity as $t\uparrow1$.  Normalize any sequence of points with
unbounded curvature.  Noncollapsing and the derivative bounds give a smooth
limit by \ref{thm:hamilton-compactness-generalized-flows}; the limit is a
noncompact $\kappa$-solution, hence cylindrical away from a possible cap.
Thus all sufficiently high-curvature points have either a cap neighborhood
or an evolving neck which extends backward for normalized time
$1+\epsilon$.

Choose $\theta<1$ so that this covers $t\ge\theta$.  On
$[0,\theta]$, asymptotic cylindricality is uniform outside one compact
spacetime set, giving necks which reach time zero.  On the remaining compact
set, positivity of curvature and smoothness give uniform curvature ratios,
diameter, volume, and derivative bounds at the curvature scale.  Enlarging
the fixed core of \ref{def:standard-surgery-cap-model} and then
$C_{\mathrm{std}}$
makes this set a controlled cap.  Smooth pointed convergence preserves each
strict estimate after a harmless enlargement of the constants, proving the
last assertion.
\end{proof}

\section{Stage 4: Continuation of Controlled Ricci Flow with Surgery}

\subsection{First-Failure Extension}

\begin{lemma}[First canonical-neighborhood failure]
\label{lem:canonical-neighborhood-first-failure}
\dcref{MT:ch17:17.2}
\source{morgan-tian}
\notready
Fix a finite induction stage and a sequence $r_a\downarrow0$.  Suppose that
for every $a$ there are controlled surgery flows with surgery controls
tending to zero for which the strong $(C,\epsilon)$-canonical-neighborhood
assumption with parameter $r_a$ fails before the stage endpoint.  For each
$a$, after passing sufficiently far down that sequence, there is a first
failure time $t_{a,b}$ and a point $x_{a,b}$ at that time such that
$R(x_{a,b})\ge r_a^{-2}$, the assumption holds before $t_{a,b}$, and no
strong $(C,\epsilon)$-canonical neighborhood exists at $x_{a,b}$.
\end{lemma}

\begin{proof}
\uses{def:controlled-ricci-flow-with-surgery,prop:evolving-neck-gluing,lem:canonical-neighborhood-stability,thm:standard-solution-canonical-neighborhoods,lem:local-surgery-standard-cap-convergence}
For fixed $a,b$, take failure points whose times decrease to the infimum of
all failure times.  The preceding time slice is compact, so after flowing
the points to a common nearby time and taking a subsequence they converge to
a point $x_{a,b}$.  Curvature continuity gives
$R(x_{a,b})\ge r_a^{-2}$.  Caps, controlled components, and round components
are open under smooth time convergence by
\ref{lem:canonical-neighborhood-stability}; if the limiting point had one of
these neighborhoods, the approximating failure points would have the same
alternative.

Suppose instead that the limiting point is the center of a strong neck.
If its backward cylinder remains in the continuing region, the spacetime
part of \ref{lem:canonical-neighborhood-stability} again transfers the neck
to the approximating points.  Otherwise that cylinder meets a newly inserted
cap.  Rescale by the surgery scale.  The local construction and
\ref{lem:local-surgery-standard-cap-convergence} identify a larger and
larger parabolic ball with the standard solution.  The alternatives in
\ref{thm:standard-solution-canonical-neighborhoods} give either a cap around
the failure point or an evolving neck overlapping the original surgery
neck; the evolving-neck gluing result \ref{prop:evolving-neck-gluing} glues the latter to
a strong neck.  Both alternatives contradict failure.  Hence the infimum is
attained at a genuine failure point with the stated preceding control.
\end{proof}

\begin{lemma}[A bounded-curvature cap encounter is standard]
\label{lem:standard-cap-encounter-control}
\dcref{MT:ch17:17.4,MT:ch17:17.5,MT:ch17:17.8}
\source{morgan-tian}
\notready
Consider rescaled first-failure flows and a backward flow line, starting in
a fixed-radius zero-time ball, which reaches a surgery cap after rescaled
time $s_a\le s'<\infty$.  If scalar curvature along the line is bounded by
$D'<\infty$ and $\bar h_a$ is the rescaled surgery scale, then for all large
$a$ there are constants $c=c(D',s')>0$, $A<\infty$, and $\theta<1$ such that
\[
 \bar h_a^2\ge c,
 \qquad s_a<\theta\bar h_a^2,
\]
and the point where the line reaches the surgery slice lies in the image of
the standard-solution ball $B(p_0,A)$ used to model the cap.
\end{lemma}

\begin{proof}
\uses{def:controlled-surgery-parameters,thm:local-neck-surgery,lem:local-surgery-standard-cap-convergence,thm:standard-solution-canonical-neighborhoods}
At the cap boundary, the surgery construction compares scalar curvature to
$\bar h_a^{-2}$ within a fixed factor $D$.  The assumed bound $D'$ along the
flow line therefore gives
\[
 \bar h_a^2\ge(2D'D)^{-1}.
\]
To reach the inserted cap, the flow line first crosses the retained collar
of the cutting neck.  In the neck coordinates of
\ref{thm:local-neck-surgery}, that collar has fixed normalized axial length,
and the local surgery construction places its inner boundary a fixed
multiple of $\bar h_a$ from the cap tip.  Thus the entry point is at distance
at most $C\bar h_a$ from the tip.  After division by $\bar h_a$, it lies in
one fixed standard ball $B(p_0,A)$.

It remains to bound elapsed time away from the standard singular time.  If
$s_a\le\bar h_a^2/2$, any $\theta\in(1/2,1)$ works.  Otherwise the curvature
bound along the line gives a fixed upper bound after cap-scale rescaling.
Choose $\theta<1$ so close to one that the scalar curvature of the standard
solution on the corresponding cap ball is larger than that bound.  Smooth
cap convergence then contradicts $s_a\ge\theta\bar h_a^2$.  This proves all
three asserted controls.
\end{proof}

\begin{lemma}[A cap-reaching flow line forces a canonical neighborhood]
\label{lem:bounded-flow-line-cap-exclusion}
\dcref{MT:ch17:17.7}
\source{morgan-tian}
\notready
In a rescaled first-failure sequence, suppose a point in a fixed-radius
zero-time ball has a backward flow line of uniformly bounded scalar
curvature which reaches a surgery cap in uniformly bounded time.  Then the
base first-failure points have strong $(C,\epsilon)$-canonical neighborhoods
for all sufficiently large indices.
\end{lemma}

\begin{proof}
\uses{lem:standard-cap-encounter-control,lem:local-surgery-standard-cap-convergence,thm:standard-solution-canonical-neighborhoods,lem:canonical-neighborhood-stability,prop:evolving-neck-gluing}
By \ref{lem:standard-cap-encounter-control}, the cap scale is bounded below,
the elapsed cap-scale time is at most some $\theta<1$, and the relevant
backward point lies in one fixed standard-solution ball.  Enlarge that ball
so that forward cap evolution contains the basepoint at time zero.  Smooth
convergence from \ref{lem:local-surgery-standard-cap-convergence} then
identifies this entire parabolic region with the corresponding standard
solution region.

Apply \ref{thm:standard-solution-canonical-neighborhoods} at the limiting
basepoint.  A cap alternative transfers directly by
\ref{lem:canonical-neighborhood-stability}.  In an evolving-neck
alternative, the neck reaches the continuing side of the original surgery
neck; their overlap charts agree after a small axial adjustment, and
\ref{prop:evolving-neck-gluing} glues them into a strong
$\epsilon$-neck.  Thus either alternative gives a forbidden canonical
neighborhood at the first-failure point.
\end{proof}

\begin{lemma}[High-curvature control in a rescaled failure slice]
\label{lem:rescaled-high-curvature-canonical-control}
\dcref{MT:ch17:17.6}
\source{morgan-tian}
\notready
After rescaling first-failure points to scalar curvature one and taking the
surgery controls to zero, every point of the zero-time slice with scalar
curvature greater than one has a strong
$(2C,2\epsilon)$-canonical neighborhood for all sufficiently large indices.
\end{lemma}

\begin{proof}
\uses{lem:canonical-neighborhood-first-failure,lem:canonical-neighborhood-stability,lem:local-surgery-standard-cap-convergence,thm:standard-solution-canonical-neighborhoods}
A nonexposed point is approached along its backward flow line by points at
times strictly before the first failure.  At those points the original
$(C,\epsilon)$-canonical-neighborhood assumption applies because their
scalar curvature is above the rescaled threshold.  Smooth convergence to
time zero and Part~(3) of
\ref{lem:canonical-neighborhood-stability} transfer a
$(2C,2\epsilon)$ neighborhood.

An exposed point lies in a newly inserted cap.  As surgery control tends to
zero, \ref{lem:local-surgery-standard-cap-convergence} identifies its
bounded cap-scale neighborhood with the standard solution, and
\ref{thm:standard-solution-canonical-neighborhoods} supplies a cap or strong
neck alternative.  Stability transfers it with doubled constants.  These
cases exhaust the zero-time slice.
\end{proof}

\begin{lemma}[Canonical transfer on a surviving first-failure slice]
\label{lem:first-failure-zero-slice-canonical-transfer}
\dcref{MT:ch17:17.1,MT:ch17:17.2,MT:ch17:17.6,MT:ch17:17.7}
\source{morgan-tian}
\uses{def:controlled-ricci-flow-with-surgery,
  lem:canonical-neighborhood-first-failure,
  lem:canonical-neighborhood-stability,lem:local-surgery-standard-cap-convergence,
  thm:standard-solution-canonical-neighborhoods,prop:evolving-neck-gluing,
  lem:bounded-flow-line-cap-exclusion}
\notready
Consider a diagonal first-failure sequence of controlled flows, with failure
points $(x_j,t_j)$ and scales
$Q_j=R_{g_j(t_j)}(x_j)\to\infty$.  Choose the threshold parameters so that
$r_j^{-2}/Q_j\le\vartheta<1$ (enlarging $r_j$ slightly before taking the
diagonal if necessary), and translate and rescale by $Q_j$ so that the
failure points have scalar curvature one at time zero.  Assume that, for every
fixed $A<\infty$, the zero-time $A$-ball has a compatible unscathed backward
parabolic neighborhood on a common positive time interval.  Then every
sequence of points in these zero-time balls whose rescaled scalar curvature is
at least one has, after passing to a subsequence, a strong
$(2C,2\epsilon)$-canonical neighborhood.  In particular the conclusion is
independent of whether the containing component has positive sectional
curvature; no positivity-to-model inference is being made.
\end{lemma}

\begin{proof}
This is the zero-slice claim in Morgan--Tian's proof of Proposition~17.1
(the source proof is \texttt{surgery.tex}, lines~4311--4355).  Let $y_j$ be
one of the stated points.  If it is exposed, the local surgery convergence
identifies a fixed cap-scale neighborhood with the standard solution, and
\ref{thm:standard-solution-canonical-neighborhoods} followed by
\ref{lem:canonical-neighborhood-stability} gives a cap alternative.

Otherwise follow $y_j$ a sufficiently short time backwards inside the assumed
compatible neighborhood.  Since the rescaled threshold is at most
$\vartheta<1$ and $R(y_j,0)\ge1$, the points at negative times sufficiently
close to zero have scalar curvature above the pre-failure threshold.  The
first-failure construction therefore gives them strong $(C,\epsilon)$-
canonical neighborhoods.  Pass to a subsequence on which the alternatives
are all of one type.  A cap, $C$-component, or round component persists by the
strict inequalities and the first part of
\ref{lem:canonical-neighborhood-stability}; a strong neck persists by the
spacetime part of that lemma.  If the backward neck reaches a surgery boundary,
\ref{lem:local-surgery-standard-cap-convergence} and the standard-solution
alternatives give a cap or an evolving neck, and
\ref{prop:evolving-neck-gluing} transfers the latter.  The bounded-flow-line
cap alternative is exactly \ref{lem:bounded-flow-line-cap-exclusion}.
Thus every stated zero-time point has the claimed doubled canonical
neighborhood.  The argument uses the canonical alternatives already present
before first failure and does not assert that positive sectional curvature by
itself supplies one of them.
\end{proof}

\begin{lemma}[First-failure noncollapsing bridge]
\label{lem:first-failure-noncollapsing-bridge}
\dcref{MT:ch16:16.1,MT:ch17:17.1,KL:II.5:II.5.2}
\source{kleiner-lott}
\source{morgan-tian}
\uses{def:surgery-flow-noncollapsing,def:ricci-flow-with-cutoff,
  def:stagewise-delta-selector,
  lem:kleiner-lott-noncollapsing-estimate}
\notready
Let $0<r_+\le r_-\le\epsilon$, $\kappa_->0$, and
$0<E_-\le E<\infty$ be fixed.  Let $a<b<c$ satisfy
$b-a\ge E_-$ and $c-a\le E$.  Consider a sequence of flows with
first-failure points $(x_j,t_j)$ with $t_j\in[b,c)$, put
$Q_j=R_{g_j(t_j)}(x_j)$, and define
the translated parabolic rescalings
\[
 \widetilde g_j(s)=Q_jg_j(t_j+Q_j^{-1}s),
 \qquad R_{\widetilde g_j(0)}(x_j)=1.
\]
Assume the original flows are three-dimensional Ricci flows with
$(r,\delta)$-cutoff, normalized compact initial data, and global stage
parameter $r\le\epsilon$, with $r\ge r_-$ on
$[a,b)$ and $r\ge r_+$ on $[b,c)$, ordinary $\kappa_-$-noncollapsing on
$[a,b)$ at every scale below $\epsilon$, and cutoff
$\delta(t)\le\bar\delta(r_-,r_+,\kappa_-,E_-,E)$ on $[a,c)$.  Assume that
the strong $r_+$ canonical-neighborhood hypothesis holds on $[b,c)$.  Let
$\kappa_+$ be the
constant supplied by
\ref{lem:kleiner-lott-noncollapsing-estimate}.  Then, after passing to a
subsequence, one may use the single constant $\kappa_*:=\kappa_+$ on every
fixed zero-time based ball: whenever the backward parabolic ball of scale
$\rho\le\epsilon$ exists, is unscathed, and has curvature at most
$\rho^{-2}$, its time-zero ball satisfies
\[
 \operatorname{Vol}_{\widetilde g_j(0)}
 B_{\widetilde g_j(0)}(x_j,\rho)
 \ge \kappa_*\rho^3
 \quad\text{or}\quad
 x_j\text{ lies in a positive-sectional-curvature component}.
\]
If an independent hypothesis excludes that positive-component alternative at
the selected basepoints, then the displayed inequality gives the fixed volume
$v_0=\kappa_*\eta^3$ on every common backward neighborhood of radius
$\eta>0$; the bridge itself does not supply that exclusion.
\end{lemma}

\begin{proof}
\uses{lem:kleiner-lott-noncollapsing-estimate}
For each compact interval $[b,c')$ before first failure, the displayed
parameter inequalities are exactly the hypotheses of
\ref{lem:kleiner-lott-noncollapsing-estimate}.  Its constant
$\kappa_+$ depends only on $(r_-,\kappa_-,E_-,E)$ and is independent of
$c'$, the particular flow, and the later scale $r_+$.  The resulting source
disjunction is preserved on the half-open first-failure interval and on each
fixed scale.  Parabolic rescaling preserves the scale-invariant inequality, so
a diagonal choice over fixed radii uses the single constant
$\kappa_*:=\kappa_+$.  The ordinary branch is recorded only when a later
argument supplies the stated nonpositive-component hypothesis; no
positive-curvature conclusion is inferred from positivity alone, and no
later-interval conclusion of the conditional noncollapsing theorem is invoked.
\end{proof}

\begin{lemma}[Pre-failure KAPPA-limit noncollapsing]
\label{lem:pre-failure-kappa-limit}
\dcref{MT:ch17:17.1,MT:ch14:14.14,MT:ch14:14.15}
\source{morgan-tian}
\uses{def:surgery-flow-noncollapsing,def:controlled-ricci-flow-with-surgery,
  def:stagewise-delta-selector,
  thm:conditional-noncollapsing-extension,
  def:ricci-flow-with-cutoff,lem:kleiner-lott-noncollapsing-estimate,
  lem:canonical-neighborhood-first-failure,
  lem:noncollapsing-predicate-pointed-limit}
\notready
Take the diagonal first-failure sequence with stage parameters
$r_a\downarrow0$, surgery controls $0<\delta_a\le\delta(r_a)$, where
$\delta(r_a)$ denotes the fixed selector value from
\ref{def:stagewise-delta-selector}, and
$\delta_a\to0$, and failure points $(x_a,t_a)$ chosen after the preceding
canonical-neighborhood interval.  After translating $t_a$ to zero and
rescaling by $Q_a=R(x_a,t_a)\to\infty$, the negative-time truncations satisfy
the Morgan--Tian/Kleiner--Lott noncollapsing predicate with one constant
$\kappa_0>0$ on every scale below $\epsilon$: every admissible test ball has
volume at least $\kappa_0\rho^3$, unless its centre lies in a compact
positive-sectional-curvature component.  The predicate is uniform on every
fixed compact spacetime region and is inherited by smooth pointed limits.
If the selected basepoint component is not positive, the ordinary volume
branch holds at every admissible based ball, including the fixed radius used
for Hamilton compactness.
\end{lemma}

\begin{proof}
This is Morgan--Tian Proposition~17.1's invocation of the KAPPA-limit
Proposition (the source statements are in \texttt{surgery.tex},
lines~2147--2165, and its first-failure application at lines~4277--4309).
Before $t_a$ the strong canonical-neighborhood assumption has the parameter
$r_a$, so the source Proposition applies to the truncated flow; the
inequalities $\delta_a\le\delta(r_a)$ and $Q_a\to\infty$ are exactly its
small-control and rescaling hypotheses.  Scaling and time translation leave
the disjunctive volume/positive-component predicate invariant.  The
first-failure stability lemma supplies the unscathed compact regions on which
the source disjunction is closed.  The positive-component alternative is retained
explicitly; an ordinary volume bound is asserted only at a centre for which
that alternative has separately been excluded.
\end{proof}

\begin{lemma}[Short backward curvature control]
\label{lem:short-backward-curvature-control}
\dcref{MT:ch17:17.9}
\source{morgan-tian}
\uses{def:stagewise-delta-selector,lem:canonical-neighborhood-first-failure,
  lem:pre-failure-kappa-limit,thm:conditional-noncollapsing-extension}
\notready
Consider the diagonal first-failure sequence with $r_a\downarrow0$,
$\delta_a\le\delta(r_a)$ and $\delta_a\to0$, where $\delta(r_a)$ is the fixed
selector value from \ref{def:stagewise-delta-selector}, with the strong
$(C,\epsilon)$-canonical-neighborhood assumption on every pre-failure slice.
Write $\widehat x_a$ for the unrescaled first-failure point and
$Q_a=R_{G_a}(\widehat x_a)\to\infty$, and assume the diagonal has been chosen
with a fixed $\vartheta<1$ satisfying
$r_a^{-2}/Q_a\le\vartheta$.
Assume the source KAPPA-limit predicate of
\ref{lem:pre-failure-kappa-limit} on the negative-time truncations.  For
every $A<\infty$ there are $D(A)<\infty$ and $\sigma(A)>0$ such that,
writing $x_a$ for the point corresponding to $\widehat x_a$ after rescaling by
$Q_a$, every sufficiently advanced rescaled flow has scalar curvature at most
$D(A)$ along each
backward flow line starting in $B(x_a,0,A)$ for as long as that line exists
in $[-\sigma(A),0]$.
\end{lemma}

\begin{proof}
\uses{thm:conditional-noncollapsing-extension,lem:pre-failure-kappa-limit,
  lem:first-failure-noncollapsing-bridge,
  lem:positive-component-persistence,
  thm:bounded-curvature-at-bounded-distance,lem:rescaled-high-curvature-canonical-control,
  thm:local-ricci-flow-derivative-estimates}
This is Claim~17.9 of Morgan--Tian (the source proof is
\texttt{surgery.tex}, lines~4541--4576).  The KAPPA-limit predicate
supplies the volume-or-positive-component disjunction on the negative-time
truncation.  If the exceptional branch occurs at a point on a backward flow
line, \ref{lem:positive-component-persistence} transports that branch to the
terminal component; otherwise the ordinary volume bound is available.  In
either case, \ref{lem:rescaled-high-curvature-canonical-control} supplies
canonical control at every unrescaled point above the threshold $Q_a$; the
ratio hypothesis ensures that these points lie in the pre-failure canonical
region supplied by \ref{lem:canonical-neighborhood-first-failure}.
Apply \ref{thm:bounded-curvature-at-bounded-distance} before rescaling, at
$\widehat x_a$ (so its hypothesis $R(\widehat x_a)\ge D_0$ holds for all
large $a$).  It gives $R_{G_a}\le D(A)Q_a$ on the unrescaled ball of radius
$A Q_a^{-1/2}$; after rescaling this is the bound $R\le D(A)$ on the
zero-time $A$-ball.  The local derivative estimate is applied only on the
unscathed parabolic neighborhoods supplied by that predicate; integrating
$|\partial_tR|\le cR^2$ along each surviving flow line gives a common short
time $\sigma(A)$ and, after pinching, a full curvature bound.  This is the
source's first-failure estimate with the diagonal and unscathedness
hypotheses made explicit.
\end{proof}

\begin{lemma}[Backward parabolic-neighborhood dichotomy]
\label{lem:backward-parabolic-neighborhood-dichotomy}
\dcref{MT:ch17:17.10}
\source{morgan-tian}
\notready
After passing to a subsequence of rescaled first-failure flows, write $x_a$
for the normalized first-failure basepoint in the $a$-th flow.  At least one
of the following alternatives occurs (the statement does not assert that the
two alternatives are logically disjoint):
\begin{enumerate}
\item for every $A<\infty$ there are $D(A),s(A)>0$ and an index $a(A)$ such
      that, for every $a\ge a(A)$, the compatible parabolic neighborhood
      $P(x_a,0,A,-s(A))$ exists and has
      $|\operatorname{Rm}|\le D(A)$;
\item on an infinite subsequence (renamed), every sufficiently late basepoint
      has a strong $(C,\epsilon)$-canonical neighborhood.
\end{enumerate}
\end{lemma}

\begin{proof}
\uses{lem:short-backward-curvature-control,lem:bounded-flow-line-cap-exclusion}
Fix an integer radius $A$.  The bound in
\ref{lem:short-backward-curvature-control} controls every backward flow line
from the zero-time $A$-ball for a uniform short time.  If all such lines
survive, their union is the required compatible parabolic neighborhood.  If
one does not survive, it reaches a surgery boundary with uniformly bounded
curvature and in uniformly bounded time, so
\ref{lem:bounded-flow-line-cap-exclusion} gives a canonical neighborhood at
the basepoint.  Diagonalize over integer $A$.  If the second outcome occurs,
retain that subsequence; otherwise the first outcome holds for every radius.
In the first-failure application below, the second outcome is excluded by the
choice of the basepoints, so the first outcome is the one used there.
\end{proof}

\begin{lemma}[Bounded curvature of a first-failure limit slice]
\label{lem:first-failure-limit-slice-bounded-curvature}
\dcref{MT:ch11:11.7}
\source{morgan-tian}
\notready
Let a time slice of a pointed first-failure limit be complete, nonflat, and
nonnegatively curved.  Suppose it is ordinarily $\kappa$-noncollapsed on all
finite scales and is the smooth limit of slices on which points above a fixed
scalar-curvature threshold have uniform canonical neighborhoods.  Then its
curvature is bounded on the whole slice.
\end{lemma}

\begin{proof}
\uses{def:surgery-flow-noncollapsing,thm:curvature-null-splitting,lem:rescaled-high-curvature-canonical-control,prop:positive-curvature-neck-scale-lower-bound}
If the curvature operator has a null eigenvalue, the strong maximum principle
and the local clause of \ref{thm:curvature-null-splitting}, applied on every
bounded parabolic ball before passing to the exhaustion, give a compatible
local product by a positively curved surface and a one-manifold.  These
products globalize on the universal cover, by the complete-slice clause of
that theorem, to $\Sigma\times\mathbb R$.  If scalar curvature never exceeds the canonical-
neighborhood threshold, it is already bounded.  Otherwise choose a point
above that threshold and lift it to the product cover.  Its lifted canonical
neighborhood is a neck or contains a
neck; whole compact-component alternatives cannot cover
$\Sigma\times\mathbb R$.  On this sufficiently thin neck the unique tangent
Ricci-null line is uniformly close to the axial line: its orthogonal planes
carry the only positive curvature eigenvalue.  The central neck sphere is
therefore transverse to the parallel line foliation and projects as a local
covering onto the surface factor.  Its image is both open and compact, hence
all of connected $\Sigma$.  Thus $\Sigma$ is compact, its continuous
curvature is bounded, and so is the curvature of the product slice and its
quotient.
If the slice is compact, boundedness follows directly by continuity.

It remains to consider a noncompact slice with positive sectional curvature.
Suppose its curvature were unbounded.  Choose points escaping compact sets
with scalar curvature tending to infinity.  The smooth-limit form of
\ref{lem:rescaled-high-curvature-canonical-control} gives a doubled canonical
neighborhood at every sufficiently high-curvature point.  A round-component
or $C$-component alternative is impossible because the component is
noncompact.  A cap alternative contains, at a scale comparable to its core
scale, the neck forming its noncompact end.  Thus every selected canonical
neighborhood contains a $2\epsilon$-neck whose scale is comparable to
$R^{-1/2}$ and hence tends to zero.  This contradicts
\ref{prop:positive-curvature-neck-scale-lower-bound}.  Thus the curvature of
the complete slice is globally bounded.
\end{proof}

\begin{lemma}[Uniform backward lifetime]
\label{lem:uniform-backward-lifetime}
\dcref{MT:ch17:17.11}
\source{morgan-tian}
\uses{lem:backward-parabolic-neighborhood-dichotomy,lem:first-failure-limit-slice-bounded-curvature,lem:short-backward-curvature-control,lem:rescaled-high-curvature-canonical-control,thm:hamilton-compactness-generalized-flows,lem:bounded-flow-line-cap-exclusion}
\notready
Suppose the first alternative of
\ref{lem:backward-parabolic-neighborhood-dichotomy} holds and the basepoints
remain first-failure points.  The backward lifetimes $s(A)$ may then be
chosen bounded below by one number $s_*>0$ independent of $A$.  Extending
maximally, either the compatible neighborhoods exist for every backward
time, or a flow line starting in one fixed-radius zero-time based ball
reaches a surgery cap at the finite maximal backward time with curvature
bounded in the original first-failure normalization.  In the infinite-
extension alternative there are, in addition, numbers $\eta_*>0$ and
$K_*<\infty$, independent of the sequence index, such that
$P(x_a,0,\eta_*,-\eta_*^2)$ is compatible and unscathed and
$|\operatorname{Rm}|\le K_*$ on it.  The radius may be decreased so that
$\eta_*\le K_*^{-1/2}$.
\end{lemma}

\begin{proof}
Apply the first alternative of
\ref{lem:backward-parabolic-neighborhood-dichotomy} at each first-failure
basepoint.  The short-backward curvature estimate gives a radius-independent
bound on a fixed fraction of the normalized ball; the rescaled canonical
control and the limit-slice bounded-curvature lemma then give a common lower
bound $s_*>0$ for the existence time of the pulled-back flow.  Hamilton
compactness supplies compatible embeddings on overlaps, so the neighborhoods
extend maximally by the following overlap-and-diagonal construction.  Choose
radii $A_m\nearrow\infty$ and, on each overlap, use the curvature and
derivative bounds from the preceding paragraph to take a subsequence whose
pulled-back metrics and time-vector fields converge on
$B(x_a,0,A_m)\times[-s_*,0]$.  A diagonal subsequence makes these limits
compatible on all overlaps, and the compatible embeddings define the maximal
backward neighborhood.  If its endpoint were finite, the same construction
on one more overlap would extend it unless a flow line left every fixed chart
or met a surgery boundary; the bounded-distance estimate rules out the
former, leaving exactly the cap encounter considered next.

If the maximal backward time were finite, the curvature bounds and the
time-vector-field coordinates would keep one flow line in a fixed-radius
zero-time ball.  The only possible obstruction is that this line meets a
surgery cap.  Its curvature remains bounded in the original normalization by
the short-backward and canonical estimates, contradicting
\ref{lem:bounded-flow-line-cap-exclusion}.  Thus the only alternatives are
the asserted infinite extension or the explicitly described cap encounter.
  The uniformity in $A$ is exactly the common-constant conclusion of Morgan--Tian
17.11.  On a fixed-radius overlap the compactness construction supplies a
uniform curvature bound $K_*$.  Shrinking the radius below
$\min\{\eta_*,K_*^{-1/2}\}$ preserves compatibility and unscathedness and
gives the final quantitative clause.
\end{proof}

\begin{theorem}[Extension of strong canonical neighborhoods]
\label{thm:strong-canonical-neighborhood-extension}
\dcref{MT:ch17:17.1}
\source{morgan-tian}
\uses{def:stagewise-delta-selector,
  thm:conditional-noncollapsing-extension}
\notready
Fix finite stage-$i$ data with ordinary $\kappa_i$-noncollapsing before
$T_i$, the stage time-gap bounds, positive pinching, and the cutoff controls
used in \ref{thm:conditional-noncollapsing-extension}; importantly, do not
assume its later-interval canonical-neighborhood conclusion.  Then there are
$0<r_{i+1}\le r_i$ and
$0<\delta_{i+1}\le\mathfrak d_i(r_{i+1})$ such that the following holds.  If a
controlled Ricci flow with surgery on $[0,T)$,
$T\in(T_i,T_{i+1}]$, has $\bar\delta\le\delta_{i+1}$ on
$[T_{i-1},T]$ and satisfies the stage-$i$ pinching, noncollapsing, and
canonical-neighborhood hypotheses before $T_i$, then it satisfies the strong
$(C,\epsilon)$-canonical-neighborhood assumption with parameter
$r_{i+1}$ throughout $[0,T)$.
\end{theorem}

\begin{proof}
\uses{def:strong-canonical-neighborhood-control,def:controlled-ricci-flow-with-surgery,def:curvature-pinched-toward-positive,thm:conditional-noncollapsing-extension,lem:canonical-neighborhood-first-failure,lem:rescaled-high-curvature-canonical-control,lem:first-failure-noncollapsing-bridge,lem:pre-failure-kappa-limit,lem:backward-parabolic-neighborhood-dichotomy,lem:first-failure-zero-slice-canonical-transfer,lem:uniform-backward-lifetime,thm:pointed-controlled-flow-compactness,thm:kappa-solution-canonical-neighborhoods,lem:bounded-flow-line-cap-exclusion,lem:positive-component-persistence}
If no pair $r_{i+1},\delta_{i+1}$ worked, choose
$r_a\downarrow0$ and, for each $a$, surgery controls
$\delta_{a,b}\downarrow0$ with counterexample flows.  Replace each threshold
parameter by a slightly larger one before diagonalizing; this preserves the
failure and gives a uniform strict ratio
$r_a^{-2}/Q_a\le\vartheta<1$, where
$Q_a=R_{g_a(t_a)}(x_a)$.  Choose the first-failure point $(x_a,t_a)$, translate
$t_a$ to zero, and rescale by $Q_a$.  First-failure minimality says exactly
that the new canonical-neighborhood assumption holds on the preceding
half-open interval; it is not being assumed on the later interval.

Apply the backward parabolic-neighborhood dichotomy.  If its second
alternative occurs, the basepoint is already canonical, contradicting its
choice.  In the first alternative, split according to the source
noncollapsing disjunction at the basepoint.  If the component is positive,
the compatible neighborhoods and the strict threshold ratio let
\ref{lem:first-failure-zero-slice-canonical-transfer} transfer a doubled
canonical neighborhood to the basepoint; the positivity persistence lemma
\ref{lem:positive-component-persistence} ensures that this branch is not being
silently replaced by an ordinary volume estimate.  If the component is not
positive, \ref{lem:first-failure-noncollapsing-bridge} supplies a fixed
ordinary volume bound on a common radius.  The uniform lifetime lemma then
provides the compatible unscathed neighborhoods required by
\ref{thm:pointed-controlled-flow-compactness}; its proof uses
\ref{lem:pre-failure-kappa-limit} to pass the source disjunction to the
negative-time limit.  If the positive branch survives in that limit, the same
zero-slice transfer gives a canonical neighborhood at the basepoint.  If it
does not survive, the limit is a genuine \(\kappa\)-solution and
\ref{thm:kappa-solution-canonical-neighborhoods} transfers a canonical
neighborhood back to the first-failure points.  A finite backward lifetime
would instead be a bounded-curvature cap encounter, excluded by
\ref{lem:bounded-flow-line-cap-exclusion}.  Every case contradicts the
chosen first failure, so some $r_{i+1}$ and $\delta_{i+1}$ satisfy the
statement.  This is Morgan--Tian Proposition~17.1 with its two source
branches and the positivity caveat kept explicit.
\end{proof}

\subsection{Surgery Continuation and Finiteness}

\begin{lemma}[Length-metric gluing for finitely many surgery pieces]
\label{lem:length-metric-gluing}
\dcref{MT:ch13:13.2}
\source{morgan-tian}
\notready
Let a smooth Riemannian manifold $X$ be cut along finitely many disjoint
embedded two-spheres into compact codimension-zero pieces $X_1,\ldots,X_m$,
with the induced length metrics and the original collar metric on each common
boundary.  Let $(Y,d_Y)$ be a metric space and let
$F_i\colon X_i\to Y$ be maps which agree on every common boundary sphere.
If each $F_i$ is $L$-Lipschitz for its intrinsic length metric, then the
glued map $F\colon X\to Y$ is $L$-Lipschitz for the ambient Riemannian
distance on $X$.
\end{lemma}

\begin{proof}
For $x,y\in X$ and $\zeta>0$, choose a piecewise smooth path from $x$ to
$y$ of length at most $d_X(x,y)+\zeta$.  The finite sphere cut can be
crossed only finitely many times after replacing the path inside collars by
radial segments, so subdivide it into subpaths contained in individual
pieces.  On each subpath the image length is at most $L$ times its source
length.  Boundary agreement makes the image concatenation continuous, and
therefore
\[
 d_Y(F(x),F(y))\le L\bigl(d_X(x,y)+\zeta\bigr).
\]
Letting $\zeta\downarrow0$ proves the claim.  This is the elementary
length-space step used by the local surgery construction; no global
Lipschitz conclusion is being inferred from boundary agreement alone.
\end{proof}

\begin{lemma}[Two-sided projective-plane exclusion survives ball surgery]
\label{lem:rp2-exclusion-under-ball-surgery}
\dcref{MT:ch15:15.8}
\source{morgan-tian}
\uses{def:controlled-ricci-flow-with-surgery,thm:local-neck-surgery}
\notready
Let $M^-$ be a closed smooth three-manifold (possibly nonorientable)
containing no embedded $\mathbb{RP}^2$ with trivial normal bundle.  Obtain
$M^+$ by finitely many
disjoint neck surgeries, each deleting a three-ball side and gluing a smooth
three-ball cap along its spherical boundary, and by discarding whole
components.  Then $M^+$ also contains no embedded $\mathbb{RP}^2$ with trivial
normal bundle.
\end{lemma}

\begin{proof}
Put a hypothetical two-sided projective plane $F$ in general position with
respect to the finite spherical surgery boundaries.  Every intersection circle
$C$ is two-sided in $F$.  Indeed, both $F$ and every surgery sphere $S$ are
two-sided, so
\[
 \nu_C^M\cong\nu_C^F\oplus\nu_F\cong\nu_C^S\oplus\nu_S,
\]
where $S$ is the surgery sphere; the three bundles $\nu_F$, $\nu_C^S$, and
$\nu_S$ are trivial, so $w_1(\nu_C^F)=0$.  Since every two-sided simple
closed curve in $\mathbb{RP}^2$
is inessential, an innermost intersection circle bounds disks on both
surfaces.  Replacing the disk in $F$ by a small push-off of the sphere disk
preserves the embedding and its normal framing and reduces the number of
intersections.

After iterating, $F$ lies entirely in a retained pre-surgery component or in
an inserted three-ball.  A three-ball contains no embedded $\mathbb{RP}^2$;
in the retained case the isotopy gives the forbidden two-sided projective
plane in $M^-$.  Discarding whole components cannot create a new embedded
surface, so $M^+$ has no such $F$.
\end{proof}

\begin{lemma}[Controlled restart after surgery]
\label{lem:controlled-surgery-restart}
\dcref{MT:ch15:15.8,MT:ch15:15.11}
\source{morgan-tian}
\notready
Let a controlled Ricci flow reach a finite singular time $T$.  In its maximal
extension, cut every sufficiently curved horn at the central sphere of a
strong $\bar\delta(T)$-neck of scale $h(T)$, discard each horn end and each
component containing no controlled-curvature point, and insert the standard
caps.  For sufficiently small control, the post-surgery metric is smooth,
retains curvature pinching, satisfies the seven controlled-flow conditions,
and is initial data for a maximal ordinary Ricci flow.  The new flow retains
the stage noncollapsing and strong canonical-neighborhood estimates.
\end{lemma}

\begin{proof}
\uses{def:controlled-surgery-parameters,def:controlled-ricci-flow-with-surgery,def:curvature-pinched-toward-positive,def:maximal-ricci-flow-extension,thm:local-neck-surgery,lem:rp2-exclusion-under-ball-surgery,thm:conditional-noncollapsing-extension,thm:strong-canonical-neighborhood-extension}
The maximal extension \ref{def:maximal-ricci-flow-extension} contains the
controlled-curvature region at time $T$.  The horn construction selects a
locally finite family of disjoint cutting necks, and compactness of the slice
makes the family finite.  Apply \ref{thm:local-neck-surgery} at each neck.
The metrics agree on the retained halves, the inserted caps are smooth, and
the local curvature comparison preserves pinching.  Discarding whole
components cannot spoil any estimate on the components which remain.

The exposed pieces are finitely many three-balls, their boundary spheres are
the prescribed central neck spheres, and the disappearing region is covered
by the canonical necks, caps, and discarded components.  By
\ref{lem:rp2-exclusion-under-ball-surgery}, the no-two-sided-$\mathbb{RP}^2$
condition also survives the ball surgeries.  Thus the seven
conditions in \ref{def:controlled-ricci-flow-with-surgery} persist.  Short-time
existence starts an ordinary flow from the compact post-surgery metric, and
uniqueness glues it to the continuing region.  Taking its maximal extension
gives the next smooth interval.  Finally
\ref{thm:conditional-noncollapsing-extension} and
\ref{thm:strong-canonical-neighborhood-extension} restore the two quantitative
stage estimates.
\end{proof}

\begin{theorem}[Map across a separating surgery: removed-region comparison]
\label{thm:separating-surgery-map}
\dcref{MT:ch15:15.12,MT:ch13:13.2}
\source{morgan-tian}
\notready
Let $({\mathcal M},G)$ be a controlled Ricci flow with surgery satisfying
Assumptions~(1)--(7) of the Morgan--Tian controlled-flow definition.  Let
$T$ be a surgery time with control parameter
$\bar\delta(T)<\delta'_0$ and surgery scale $h(T)<R_0^{-1/2}$, where
$\delta'_0,R_0$ are the constants of
\ref{thm:local-neck-surgery}.  Fix a regular $t'<T$ sufficiently close to
$T$, let $X(t')$ be a presurgery component, and let $X(T)$ be a component
obtained from it at $T$.  If every cutting two-sphere in $X(t')$ is
separating, then the time-flow embedding has an open bounded-curvature
domain $\Omega\subset X(t')$ whose metric limit exists and whose limiting
identification is an isometry onto an open subset of $X(T)$.  The
identification extends continuously to a map
\[
 F\colon X(t')\longrightarrow X(T)
\]
such that, for every regular time $t<T$ sufficiently close to $T$, its
restriction to $X(t')\setminus\Omega$, measured with the time-$t$ pulled-back
metric, is distance nonincreasing into $X(T)$.  The radial cross-neck
extension is the one supplied by the local surgery theorem.
\end{theorem}

\begin{proof}
\uses{def:controlled-ricci-flow-with-surgery,thm:local-neck-surgery,
  lem:radial-surgery-map-nonexpanding}
This is Proposition~15.12 of Morgan--Tian with its controlled-flow,
small-control, scale, regular-time, and separating-sphere hypotheses made
explicit.  The bounded-curvature region and its limiting isometry are the
first part of that proposition.  For each separating cutting sphere there is
one side not leading to the selected persistent component; the local surgery
construction sends that side to the standard cap by the radial map of
\ref{lem:radial-surgery-map-nonexpanding}.  The local theorem identifies the
map on the retained neck and supplies the cross-neck agreement.  The source
proposition gives, simultaneously for the finite family of surgery regions,
the distance-decreasing estimate on the discarded region.  No global
Lipschitz estimate is included here; smooth convergence on compact subsets
alone does not control distances up to the blow-up frontier.
\end{proof}

\begin{lemma}[Global Lipschitz extension across trivial surgery]
\label{lem:global-surgery-map-lipschitz}
\dcref{MT:ch18:18.18,MT:ch15:15.12}
\source{morgan-tian}
\uses{thm:separating-surgery-map,def:controlled-ricci-flow-with-surgery,
  thm:local-neck-surgery,lem:radial-surgery-map-nonexpanding,
  lem:length-metric-gluing}
\notready
In the setting of
\ref{thm:separating-surgery-map}, assume in addition that the selected
post-surgery component has a decomposition
\[
 X(T)=C(T)\cup K_1(T)\cup\cdots\cup K_m(T),
\]
where $C(T)$ is a compact codimension-zero persistent core, the $K_i(T)$ are
finitely many compact surgery caps with pairwise disjoint interiors, and all
common boundaries are the cutting two-spheres.  Assume every cutting sphere
is homotopically trivial in $X(t')$.  For $t'<T$ sufficiently close to $T$,
the backward flow identifies $C(T)$ with a compact submanifold
$C(t')\subset X(t')$ by an embedding $n_{t'}$, and
\[
 n_{t'}^{*}\bigl(g(t')|_{C(t')}\bigr)\longrightarrow G(T)|_{C(T)}
 \quad\text{in }C^\infty\text{ on }C(T).
\]
Then, for every $\eta>0$, there is $t_\eta<T$ such that for every regular
$t'\in(t_\eta,T)$ the inverse core map $n_{t'}^{-1}$ extends over the
finitely many cap pieces to a map
\[
 \psi_{t'}\colon X(t')\longrightarrow X(T)
\]
which maps $X(t')\setminus C(t')$ into the union of the caps and satisfies
\[
 \operatorname{Lip}_{(X(t'),g(t'))\to(X(T),G(T))}(\psi_{t'})\le 1+\eta
\]
for that $\eta$.  This metric extension statement deliberately makes no
homotopy-type claim; the old-type versus homotopy-sphere-descendant
alternative is supplied by the separate topology contract before any
homotopy consumer invokes $\psi_{t'}$.
\end{lemma}

\begin{proof}
The compact-core/finite-cap decomposition is the surgery description used in
Morgan--Tian Chapter~18.  Homotopical triviality of each cutting sphere makes
it separating, so the inverse of the backward core embedding extends across
the corresponding cap by the radial local-surgery map.  The $C^\infty$
convergence on the compact core
gives arbitrarily small multiplicative distortion there.  Apply
\ref{thm:separating-surgery-map} to the finitely many complementary pieces;
its removed-region distance-decrease estimate controls every cap-crossing
subpath.  The explicit length-space decomposition
\ref{lem:length-metric-gluing} converts these intrinsic piece estimates to
the ambient distance.  Taking the maximum of the core distortion and the cap
estimate gives the displayed global $(1+\eta)$ bound.  This is the compact-core
corollary used in the Morgan--Tian width comparison, not a consequence of
boundary agreement alone.
\end{proof}

\begin{lemma}[Uniform volume loss at a cutoff surgery]
\label{lem:cutoff-surgery-volume-loss}
\dcref{KL:II.4:volumelossremark,MT:ch17:17.12}
\source{kleiner-lott}
\uses{def:ricci-flow-with-cutoff,lem:standard-surgery-cap-model}
\notready
Fix the dimension-three neck parameter $\epsilon$ and choose the cutoff
parameter $\bar\delta$ below the standard cap threshold.  There is a constant
$c_{\rm vol}=c_{\rm vol}(\epsilon,\bar\delta)>0$ such that, at every cutoff
surgery on a central neck of scale $h$, the total volume after inserting the
standard cap and discarding the exposed component satisfies
\[
 \operatorname{Vol}(M_t^-)-\operatorname{Vol}(M_t^+)
 \ge c_{\rm vol}h^3.
\]
The same bound holds componentwise whenever a neck is cut, and it is uniform
over all scales and all surgery times for which the displayed parameters are
used.
\end{lemma}

\begin{proof}
Use the neck coordinate and rescale by $h^{-2}$, so the central neck has
scale one.  On the fixed half-neck retained in the exposed region, the strong
$\bar\delta$-neck estimate gives $C^2$ closeness to a unit round cylinder.
Its axial interval has length at least a fixed multiple of
$\bar\delta^{-1}$, hence its volume is at least
$a_0\bar\delta^{-1}$ for a universal $a_0>0$.  The standard cap model in
\ref{lem:standard-surgery-cap-model} has a fixed rescaled volume at most
$A_0$; the $C^2$ cap comparison changes this by at most $o_{\bar\delta}(1)$.
Choose $\bar\delta$ so that
$a_0\bar\delta^{-1}-(A_0+o_{\bar\delta}(1))\ge c_{\rm vol}>0$.
The surgery region is disjoint from all other cutting regions, so the local
loss adds to the global volume loss.  Scaling back multiplies volume by
$h^3$, proving the claim.  This is the quantitative content of the
Kleiner--Lott volume-loss remark.
\end{proof}

\begin{theorem}[Finite surgery count on bounded intervals]
\label{thm:finite-interval-surgery-bound}
\dcref{MT:ch17:17.12}
\source{morgan-tian}
\uses{lem:finite-horizon-surgery-phi-envelope}
\notready
Let a controlled Ricci flow with surgery be defined on $[0,T)$ with
$T\le T_{i+1}$ and with the control functions fixed on that interval.  Assume
the positive nonincreasing cutoff function $\bar\delta$ is defined through
$T_{i+1}$ and that $r(t)\ge r_{i+1}>0$ there.  Put
\[
 \Phi:=\Phi_{\rm HI}^{T_{i+1},g(0)}
 \quad\text{from \ref{lem:finite-horizon-surgery-phi-envelope}},
 \qquad
 h_*:=H_*(\Phi,T_{i+1},r_{i+1},
              \bar\delta(T_{i+1}))>0,
\]
where $H_*$ is the bounded-stage lower envelope from
\ref{lem:ricci-flow-surgery-horn-scale}.  Then every cutting neck in $[0,T)$
has $h(t)\ge h_*$.  There is
\[
 N=N\bigl(\operatorname{Vol}(M_0,g(0)),T_{i+1},r_{i+1},
             \bar\delta(T_{i+1}),h_*\bigr)<\infty
\]
such that the flow has at most $N$ surgery times in $[0,T)$.  In particular,
surgery times cannot accumulate in a bounded interval.
\end{theorem}

\begin{proof}
\uses{def:controlled-surgery-parameters,def:controlled-ricci-flow-with-surgery,
  def:curvature-pinched-toward-positive,
  lem:finite-horizon-surgery-phi-envelope,
  lem:ricci-flow-surgery-horn-scale,
  lem:cutoff-surgery-volume-loss}
Away from surgery, curvature pinching gives $R\ge-6$ and hence
\[
 \frac d{dt}\operatorname{Vol}(M_t)
 =-\int_{M_t}R\,d\operatorname{vol}
 \le6\operatorname{Vol}(M_t).
\]
Thus the total volume immediately before every surgery is at most
\[
 e^{6T_{i+1}}\operatorname{Vol}(M_0)
\]
plus the earlier downward jumps.

A neck surgery has the uniform loss supplied by
\ref{lem:cutoff-surgery-volume-loss}.  Since $r$ and $\bar\delta$ are
positive and nonincreasing on the bounded interval,
\[
 r(t)\ge r_{i+1},\qquad
 \bar\delta(t)\ge\bar\delta(T_{i+1})>0.
\]
The definition of the controlled surgery parameters and the bounded-stage
envelope in \ref{lem:ricci-flow-surgery-horn-scale} therefore give
$h(t)\ge h_*$ at every cutting time.  Hence every neck surgery has the fixed
lower loss $v_*:=c_{\rm vol}h_*^3>0$, so
the number of cutting operations is bounded by the exponential volume bound
divided by $v_*$.

If $S$ is the number of cutting operations and $D$ the number of whole
components discarded, the number of components is at most $N_0+S-D$.
Normalization gives a fixed positive lower volume for every initial
component, hence bounds $N_0$ by the initial volume.  It follows that $D$ is
bounded as well.  Each surgery time performs at least one cutting or discards
at least one component, which proves the stated bound.
\end{proof}

\subsection{All-Time Controlled Extension}

\begin{theorem}[All-time controlled Ricci flow with surgery]
\label{thm:all-time-controlled-surgery-flow}
\dcref{MT:ch15:15.9,MT:ch15:15.10}
\source{morgan-tian}
\source{kleiner-lott}
\uses{def:ricci-flow-with-cutoff,def:stagewise-delta-selector}
\notready
There are positive nonincreasing sequences
$\mathbf K=(\kappa_i)_{i\ge0}$,
$\boldsymbol\delta=(\delta_i)_{i\ge0}$, and
$\mathbf r=(r_i)_{i\ge0}$ with the following property.  For every positive
nonincreasing function $\bar\delta\colon[0,\infty)\to(0,\infty)$ which has
no transition on $[0,T_0]$ and satisfies
\[
\bar\delta(t)<\delta_{i+1}\le\delta_0^{\rm cap}
\qquad\text{at every transition }T_i<t\le T_{i+1},
\]
every normalized compact three-manifold containing no embedded
$\mathbb{RP}^2$ with trivial normal bundle is the initial slice of a controlled
Ricci flow with surgery on $[0,\infty)$ using the Kleiner--Lott
$(r,\bar\delta)$-cutoff of \ref{def:ricci-flow-with-cutoff}.  For
$i\ge0$, $t\in[T_i,T_{i+1})$, $x\in M_t$, and
$0<\rho\le\epsilon$, every parabolic ball satisfying the curvature hypothesis in
\ref{def:surgery-flow-noncollapsing} satisfies its source disjunction with
constant $\kappa_{i+1}$: either its volume is at least
$\kappa_{i+1}\rho^3$, or $x$ lies in a component with positive sectional
curvature.  Any later use requiring a volume bound is restricted to the first
branch (or supplies a separate compact positive-component estimate).
The flow satisfies the strong $(C,\epsilon)$-canonical-neighborhood assumption
with parameter $r_{i+1}$, has curvature pinched toward positive, and uses
surgery control at most $\delta_{i+1}$.  Its surgery times are finite on every
bounded interval.
\end{theorem}

\begin{proof}
\uses{def:surgery-flow-noncollapsing,def:controlled-surgery-parameters,def:controlled-ricci-flow-with-surgery,def:curvature-pinched-toward-positive,thm:normalized-short-time-ricci-flow,thm:conditional-noncollapsing-extension,thm:strong-canonical-neighborhood-extension,lem:controlled-surgery-restart,thm:finite-interval-surgery-bound}
Apply \ref{thm:normalized-short-time-ricci-flow} to the normalized initial
metric.  It gives the ordinary flow through $T_0=2^{-5}$, its initial
noncollapsing constant, and the curvature bounds which make the stage-zero
canonical and pinching requirements automatic after fixing the universal
constants.

Choose the finite parameter sequences inductively.  Once the data through
stage $i$ are fixed, \ref{thm:conditional-noncollapsing-extension} supplies
$\kappa_{i+1}\le\kappa_i$ and the fixed selector value
$\mathfrak d_i(r_{i+1})$ for each proposed $r_{i+1}$.  Then
\ref{thm:strong-canonical-neighborhood-extension} supplies
$r_{i+1}\le r_i$ and
$\delta_{i+1}\le\mathfrak d_i(r_{i+1})$.

In the finite control sequence, replace the previous terminal entry
$\delta_i$ by $\delta_{i+1}$ and append a second copy of
$\delta_{i+1}$.  Thus $\bar\delta\le\delta_{i+1}$ throughout
$[T_{i-1},T_{i+1})$, exactly the overlap required by both extension theorems.
Earlier entries are unchanged, and $\delta_{i+1}\le\delta_i$, so the resulting
finite sequence remains positive and nonincreasing.  Each entry stabilizes
after it ceases to be terminal; the stabilized entries define
$\boldsymbol\delta$, while the appended $r_{i+1}$ and $\kappa_{i+1}$ define
$\mathbf r$ and $\mathbf K$.

Run the maximal ordinary flow from the current slice.  If it becomes singular
before $T_{i+1}$, apply \ref{lem:controlled-surgery-restart}; if no component
remains, declare every later slice empty.  Otherwise repeat from the new
compact slice.  The two extension theorems retain the quantitative hypotheses
after each restart.  By \ref{thm:finite-interval-surgery-bound}, only finitely
many restarts occur before $T_{i+1}$, so the construction reaches the stage
endpoint.

Induction over $i$ gives compatible controlled flows on all bounded time
intervals.  Their union is a Ricci flow with surgery on $[0,\infty)$ and has
every property asserted in the statement.
\end{proof}
